\documentclass[14pt,a4paper]{extarticle}

\usepackage[T2A]{fontenc}
\usepackage[utf8]{inputenc}
\usepackage[russian]{babel}
\usepackage[a4paper,margin=2cm]{geometry}
\usepackage{amsmath,amssymb}
\usepackage{graphicx}
\usepackage{booktabs}
\usepackage{longtable}
\usepackage{float}
\usepackage{array}
\usepackage{caption}
\usepackage{subcaption}
\usepackage{hyperref}
\usepackage{listings}
\usepackage{xcolor}
\usepackage{setspace}

\onehalfspacing
\graphicspath{{outputs/plots/}}

\hypersetup{
    colorlinks=true,
    linkcolor=black,
    urlcolor=blue
}

\definecolor{codebg}{RGB}{245,245,245}
\definecolor{codecomment}{RGB}{0,128,0}
\definecolor{codekeyword}{RGB}{0,0,180}
\definecolor{codestring}{RGB}{163,21,21}

\lstdefinestyle{pythonstyle}{
    language=Python,
    backgroundcolor=\color{codebg},
    basicstyle=\ttfamily\small,
    keywordstyle=\color{codekeyword}\bfseries,
    commentstyle=\color{codecomment},
    stringstyle=\color{codestring},
    numbers=left,
    numberstyle=\tiny,
    stepnumber=1,
    numbersep=8pt,
    showstringspaces=false,
    breaklines=true,
    frame=single,
    tabsize=4
}

\title{Лабораторная работа \\ Методы одномерного поиска}
\author{
    TODO: ФИО студентов \\
    TODO: Номер группы
}
\date{TODO: 2026}

\begin{document}

\begin{titlepage}
    \begin{center}
        Министерство науки и высшего образования Российской Федерации \\[0.5cm]
        TODO: Полное название университета \\[0.5cm]
        TODO: Название института / факультета / кафедры \\[3cm]

        {\large Отчет по лабораторной работе} \\[0.5cm]
        {\LARGE \textbf{Методы одномерного поиска}} \\[2cm]

        \begin{flushright}
            Выполнили: \\[0.2cm]
            TODO: ФИО, группа \\[0.5cm]
            Проверил: \\[0.2cm]
            TODO: ФИО преподавателя
        \end{flushright}

        \vfill
        TODO: Город \\[0.2cm]
        2026
    \end{center}
\end{titlepage}

\tableofcontents
\newpage

\section{Постановка задачи}

Цель работы --- изучить и сравнить методы одномерной оптимизации на различных типах целевых функций.

В рамках лабораторной работы требуется:
\begin{enumerate}
    \item Реализовать собственные методы одномерного поиска.
    \item Протестировать методы на унимодальных и многомодальной функциях.
    \item Сравнить собственные реализации с библиотечным методом Брента из \texttt{scipy.optimize.minimize\_scalar}.
    \item Построить таблицы и графики, отражающие зависимость эффективности методов от требуемой точности.
\end{enumerate}

\section{Используемые модели задач}

\subsection{``Хорошая'' унимодальная функция}

Рассматривалась функция с явно выраженным минимумом:
\[
f_1(x) = TODO
\]

Интервал поиска:
\[
[a_1, b_1] = \text{TODO}
\]

\textbf{Промежуточный вывод.} TODO: кратко описать, почему функция удобна для базового сравнения методов.

\subsection{Унимодальная функция с особенностями поведения}

Рассматривалась унимодальная функция с плато, асимметрией или иными особенностями:
\[
f_2(x) = TODO
\]

Интервал поиска:
\[
[a_2, b_2] = \text{TODO}
\]

\textbf{Промежуточный вывод.} TODO: отметить, какие особенности функции должны повлиять на работу методов.

\subsection{Многомодальная функция}

TODO: добавить после реализации многомодального сценария.

\section{Описание методов оптимизации}

\subsection{Метод пассивного поиска}

TODO: кратко описать идею метода, способ выбора сетки и критерий останова.

\subsection{Метод Брента}

В работе для сравнения использовалась библиотечная реализация метода Брента:
\begin{center}
\texttt{scipy.optimize.minimize\_scalar(..., method='brent')}
\end{center}

TODO: кратко описать преимущества метода Брента как комбинированного подхода.

\subsection{Другие методы}

TODO: после реализации добавить описания:
\begin{itemize}
    \item метода дихотомии;
    \item метода золотого сечения;
    \item метода Фибоначчи;
    \item метода парабол.
\end{itemize}

\section{Структура программной реализации}

Программа реализована на языке Python с использованием объектно-ориентированного подхода. Основные сущности:
\begin{itemize}
    \item базовый абстрактный класс целевой функции;
    \item конкретные модели тестовых функций;
    \item базовый абстрактный класс оптимизатора;
    \item классы конкретных методов поиска;
    \item класс эксперимента для серийного запуска и сбора статистики;
    \item функции вывода таблиц и построения графиков.
\end{itemize}

\textbf{Промежуточный вывод.} TODO: кратко описать, почему такая архитектура удобна для расширения командой проекта.

\section{Результаты вычислительных экспериментов}

\subsection{Таблицы результатов}

Для каждой унимодальной функции и каждого метода требуется привести таблицы зависимости количества итераций и числа вычислений функции от точности.

\begin{table}[H]
    \centering
    \caption{Пример таблицы результатов для первой функции}
    \begin{tabular}{|c|c|c|c|c|}
        \hline
        № & Метод & $\varepsilon$ & Количество итераций & Количество вычислений функции \\
        \hline
        1 & TODO & TODO & TODO & TODO \\
        \hline
        2 & TODO & TODO & TODO & TODO \\
        \hline
    \end{tabular}
\end{table}

\textbf{Промежуточный вывод.} TODO: описать, какие методы оказались более экономичными по итерациям и вызовам функции.

\subsection{Графики зависимости эффективности от точности}

\begin{figure}[H]
    \centering
    \includegraphics[width=0.85\textwidth]{TODO_first_function_metrics.png}
    \caption{Зависимость числа итераций и числа вычислений функции от точности для первой функции}
\end{figure}

\begin{figure}[H]
    \centering
    \includegraphics[width=0.85\textwidth]{TODO_second_function_metrics.png}
    \caption{Зависимость числа итераций и числа вычислений функции от точности для второй функции}
\end{figure}

\textbf{Промежуточный вывод.} TODO: прокомментировать форму графиков и различия в поведении методов.

\subsection{Графики динамики интервала неопределенности}

\begin{figure}[H]
    \centering
    \includegraphics[width=0.75\textwidth]{TODO_interval_plot_1.png}
    \caption{Динамика интервала неопределенности для выбранного метода и функции}
\end{figure}

\begin{figure}[H]
    \centering
    \includegraphics[width=0.75\textwidth]{TODO_interval_plot_2.png}
    \caption{Динамика интервала неопределенности для второго примера}
\end{figure}

\textbf{Промежуточный вывод.} TODO: отметить характер сужения интервала для разных методов.

\subsection{Сравнение с библиотечным методом Брента}

TODO: добавить сравнительный анализ собственных реализаций и метода Брента на всех функциях.

\textbf{Промежуточный вывод.} TODO: указать, в каких случаях библиотечный метод оказался устойчивее или эффективнее.

\subsection{Многомодальная функция и предварительная разведка}

TODO: после реализации многомодальной функции описать:
\begin{itemize}
    \item применение пассивного поиска для грубой оценки интервала;
    \item запуск активного метода на найденном интервале;
    \item сравнение с запуском без предварительной разведки.
\end{itemize}

\textbf{Промежуточный вывод.} TODO: сделать вывод о полезности этапа разведки.

\section{Общий вывод}

TODO: сформулировать общий вывод по лабораторной работе. Рекомендуется отразить:
\begin{itemize}
    \item какие методы оказались наиболее точными;
    \item какие методы потребовали меньше вычислений;
    \item как свойства целевой функции влияют на эффективность методов;
    \item чем библиотечный метод отличается от собственных реализаций;
    \item какие ограничения и особенности были выявлены в ходе экспериментов.
\end{itemize}

\appendix

\section{Листинг программы}

Ниже приводится основной файл программной реализации.

\lstinputlisting[style=pythonstyle]{lab1_linear_search.py}

\end{document}
